\documentclass[11pt]{article}

\usepackage[landscape,margin=0.7in]{geometry}
\usepackage{longtable}
\usepackage{booktabs}
\usepackage{array}
\usepackage{xcolor}
\usepackage{ragged2e}
\usepackage{hyperref}
\hypersetup{colorlinks=true, linkcolor=black, urlcolor=blue!60!black}

% Confidence colour coding
\newcommand{\confhigh}[1]{\textcolor{green!45!black}{\textbf{#1}}}
\newcommand{\confmed}[1]{\textcolor{blue!55!black}{#1}}
\newcommand{\conflow}[1]{\textcolor{gray}{#1}}

% Ragged-right paragraph columns (avoid ugly justification in narrow cells)
\newcolumntype{P}[1]{>{\RaggedRight\arraybackslash}p{#1}}

\title{\textbf{Research Inventory:\\ Preserved Findings for Future Papers}}
\author{Oxford University Ledger Project, 1700--1900}
\date{Companion to \texttt{paper.tex} (core leapfrogging evidence package). May 2026.}

\begin{document}

\maketitle

\noindent
\textbf{Purpose.} The Oxford ledger corpus (1{,}581 enriched pages, 1700--1900) supports several
distinct research questions beyond the central leapfrogging story developed in the main paper.
This inventory catalogues each preserved finding with its empirical support, methodological
approach, a working confidence rating, a candidate theoretical angle, and a possible future-paper
direction. The intent, following advisor guidance, is to treat the archive as a long-term research
platform rather than a single-paper project. Confidence ratings and theoretical/future-direction
columns are working judgments grounded in the current results and are expected to be revised by the
authors as the research programme matures.

\medskip
\noindent\textbf{Confidence legend:}\quad
\confhigh{HIGH} = robust across specifications;\quad
\confmed{MEDIUM} = directionally clear but with caveats;\quad
\conflow{LOW} = suggestive / specification-sensitive.

\bigskip

\footnotesize
\setlength{\extrarowheight}{3pt}
\begin{longtable}{P{2.4cm} P{4.2cm} P{3.2cm} P{1.5cm} P{4.0cm} P{4.2cm}}
\toprule
\textbf{Finding} & \textbf{Empirical support} & \textbf{Methodological approach} &
\textbf{Confidence} & \textbf{Theoretical angle} & \textbf{Future-paper direction} \\
\midrule
\endfirsthead

\toprule
\textbf{Finding} & \textbf{Empirical support} & \textbf{Methodological approach} &
\textbf{Confidence} & \textbf{Theoretical angle} & \textbf{Future-paper direction} \\
\midrule
\endhead

\midrule
\multicolumn{6}{r}{\textit{continued on next page}} \\
\endfoot

\bottomrule
\endlastfoot

% ---- 1. Capability before mission ----
Capability expansion precedes mission transformation (L3 $\rightarrow$ L4)
& Rolling 30-yr correlation shows L3 leading L4 by 5 yr in pre-1750 windows (advantage 0.5--0.7); the relationship turns synchronous after $\sim$1750. First-difference cross-correlation peak is directionally consistent but \emph{not} statistically significant; first-entry timing is threshold-sensitive.
& Cross-lag correlation with block-bootstrap CIs; distributed-lag models (levels and first differences); rolling-window correlation; state-transition entry timing.
& \confmed{MEDIUM}
& Dynamic-capabilities / organisational-learning view: capability stocks accumulate before strategic redeployment. The shift from pre-1750 lead--lag to post-1750 synchrony is itself a distinctive object of study.
& ``Capability accumulation as a precursor to mission change'': a focused lead--lag paper exploiting the time-varying structure (sequential in calm periods, simultaneous under shocks). \\

% ---- 2. Strict ordering ----
Strict L1 $\rightarrow$ L2 $\rightarrow$ L3 $\rightarrow$ L4 stage ordering
& Change-point ordering is only partial (L4 has the largest relative break, but break-\emph{year} ordering is mixed); transition-date crossings are sensitive to the chosen threshold. No robust strict ordering emerges.
& Binary-segmentation change-point detection; transition-date crossings at multiple quantile thresholds.
& \conflow{LOW}
& Critique of stage / maturity models: transformation dimensions are not traversed in a fixed universal sequence.
& ``Against stage models'': reframing four-stage frameworks as heterogeneous, shock-contingent pathways rather than deterministic sequences. \\

% ---- 3. Granger chaining ----
Granger-causal chaining across levels ($L_k$ predicts $L_{k+1}$)
& Mostly insignificant across lag specifications and pairwise tests; no evidence of predictive chaining between adjacent levels.
& Sequential Granger causality (HAC-robust); local-projection impulse responses.
& \conflow{LOW}
& Useful null result: undermines mechanistic ``each stage triggers the next'' accounts of organisational modernisation.
& Methodological / robustness companion, or a short note on the limits of Granger logic for slow-moving institutional series. \\

% ---- 4. Language modernisation ----
Semantic / language modernisation (Latin $\rightarrow$ English; vocabulary evolution)
& English-language share of entries rises from a low base across the 18th century; vocabulary diversification trends documented in standalone text analyses.
& Per-row language classification (Latin / English / mixed); vocabulary type--token ratios; text-trend analysis over the enriched corpus.
& \confmed{MEDIUM}--\confhigh{HIGH}
& Cultural-cognitive institutional legitimacy: vernacularisation as a marker of administrative and cultural modernisation.
& ``Language as a marker of institutional modernisation'': vernacularisation of record-keeping as a window on cultural change in a conservative institution. \\

% ---- 5. Income diversification ----
Income diversification dynamics (concentration $\rightarrow$ diversification)
& Income-source concentration (HHI) declines across the 19th century as new revenue categories enter.
& Annual Herfindahl index of income categories; entry/exit accounting of revenue streams.
& \confmed{MEDIUM}
& Resource-dependence and portfolio / risk-diversification theory: broadening the revenue base alters institutional autonomy.
& ``Revenue diversification and institutional resilience'': linking income-base breadth to capacity for transformation. \\

% ---- 6. Charity institutionalisation ----
Charity institutionalisation (informal alms $\rightarrow$ formal endowment)
& Charity-related categories shift from ad-hoc disbursements toward recurring, formalised payments over the 18th--19th centuries.
& Category-share decomposition; hypothesis-driven longitudinal analysis of the enriched \texttt{category} field.
& \confmed{MEDIUM}
& Institutionalisation theory (DiMaggio \& Powell): the routinisation and formalisation of previously discretionary giving.
& ``From alms to endowment'': the institutionalisation of charitable activity within a corporate body. \\

% ---- 7. Corporate-market transition ----
Corporate / market transition (in-kind \& ad-hoc $\rightarrow$ cash \& market-mediated)
& Payment composition shifts toward salaried, regular money flows; in-kind and one-off disbursements decline.
& Payment-type composition over time; supplier-network structure from the enriched \texttt{category} / \texttt{direction} / \texttt{payment\_period} fields.
& \conflow{LOW}--\confmed{MEDIUM}
& Market formation and monetisation (Polanyian ``great transformation'' at the level of a single institution).
& ``Monetisation of an institution'': the transition from in-kind, relational exchange to cash, market-mediated provisioning. \\

% ---- 8. Payment modernisation ----
Payment modernisation (irregular $\rightarrow$ periodic, salary-based)
& The payment-modernity index (L2) rises smoothly across 1700--1900 and serves as a stable transformation proxy throughout the main analysis.
& Payment-period regularity index aggregated per year; 10-yr rolling smoothing.
& \confhigh{HIGH}
& Weberian bureaucratisation / administrative rationalisation: the regularisation of payment as a core modernising process.
& Measurement companion paper, or supporting evidence integrated into the operational-axis discussion of the main paper. \\

% ---- 9. Resource slack ----
Resource slack enables leapfrogging (endowment / land-rent income)
& Rising land-rent income provides the financial capacity for capability investment during reform-shock periods; documented as an enabling mechanism in the income-composition analyses.
& Income-composition decomposition; interrupted time series on land-rent income alongside the expenditure dimensions.
& \confmed{MEDIUM}--\confhigh{HIGH}
& Organisational-slack theory (Cyert \& March): slack as the precondition that permits investment ahead of strategic certainty.
& Slack as a moderator of leapfrogging --- strongest as an integrated mechanism \emph{within} the main paper, but extendable to a comparative study of resource-constrained vs.\ well-endowed institutions. \\

\end{longtable}

\normalsize
\bigskip

\noindent\textbf{How to use this inventory.}
\begin{itemize}
\item Each row is a candidate paper or supplementary contribution; group rows by theoretical angle
      when scoping new projects.
\item Source artifacts are reproducible: the CSVs and scripts behind each finding live under
      \texttt{experiments/reports/} and \texttt{experiments/analysis/} in the project repository.
\item Treat the Confidence, Theoretical angle, and Future-paper columns as a working document;
      update them as new evidence accrues.
\item The capability-before-mission row is the most paper-ready candidate beyond the core
      leapfrogging story --- the pre-1750 lead--lag shifting to post-1750 synchrony is a
      distinctive empirical contribution in its own right.
\end{itemize}

\end{document}
